\section{Cryptographic Foundations}

\subsection{(10 điểm) Các khái niệm cơ bản}

\textbf{(a) (3 điểm) Định nghĩa ba mục tiêu an ninh chính của mật mã học bằng lời của bạn và đưa ra một ví dụ thực tế cho mỗi mục tiêu mà không được đề cập trực tiếp trong bài giảng.}

\begin{center}
\textbf{Trả lời:}
\end{center}

\textbf{Confidentiality (Bảo mật)}: Chỉ những bên được phép mới đọc được dữ liệu.
\begin{itemize}
    \item Nói cách khác, dù dữ liệu có bị chặn trên đường đi, nội dung vẫn không bị lộ.
    \item Ví dụ: Vault của trình quản lý mật khẩu (1Password/Bitwarden) được mã hóa end-to-end bằng khóa của người dùng thay vì lưu plain text trong cơ sở dữ liệu.
\end{itemize}

\textbf{Authentication (Xác thực)}: Xác minh danh tính thực sự của người gửi/người dùng.
\begin{itemize}
    \item Điều này giúp ta biết ``đúng người, đúng khóa'', tránh mạo danh.
    \item Ví dụ: Đăng nhập dùng khóa bảo mật phần cứng FIDO2/WebAuthn (YubiKey) để chứng minh danh tính mà không cần mật khẩu.
\end{itemize}

\textbf{Integrity (Toàn vẹn)}: Dữ liệu không bị thay đổi trong lưu trữ/truyền tải; nếu thay đổi sẽ bị phát hiện.
\begin{itemize}
    \item Tức là bất kỳ khác biệt nhỏ nào so với bản gốc cũng sẽ bị phát hiện một cách đáng tin cậy.
    \item Ví dụ: Kiểm tra chữ ký/tệp checksum SHA-256 kèm theo các bản cài đặt phần mềm (release checksum + signature) trước khi cài đặt.
\end{itemize}

\textbf{(b) (3 điểm) Giải thích nguyên lý Kerckhoff và lý do tại sao nguyên lý này vẫn là nền tảng của mật mã học hiện đại. Hãy đưa ra một ví dụ về một hệ thống bảo mật vi phạm nguyên lý này và mô tả những hậu quả tiềm ẩn.}

\begin{center}
\textbf{Trả lời:}
\end{center}

\textbf{Nguyên lý Kerckhoff}: Hệ mật mã vẫn phải an toàn ngay cả khi toàn bộ thuật toán/thiết kế/mã nguồn đều công khai; chỉ khóa là bí mật.

Có nghĩa là hệ mật mã chỉ cần bảo mật duy nhất khoá, các thứ khác (như là thuật toán, thiết kế, ...) thì có thể (hoặc là khuyến khích) công khai.

\textbf{Vì sao nền tảng}:

\begin{itemize}
    \item Tránh ``security through obscurity''; thuật toán lộ không làm sụp đổ hệ thống.
    \item Tập trung bảo vệ khóa (bí mật duy nhất), đơn giản hóa vận hành.
    \item Khuyến khích kiểm toán/đánh giá ngang hàng, tăng minh bạch và độ tin cậy, giảm nguy cơ backdoor.
\end{itemize}

\textbf{Ví dụ vi phạm}: Thiết bị IoT dùng ``mật mã tự chế'' đóng kín, dựa vào việc giấu thuật toán.
\begin{itemize}
    \item \textbf{Hậu quả}: Khi bị reverse-engineering rò thuật toán, kẻ tấn công có thể giải mã/giả mạo; không có đánh giá cộng đồng nên lỗ hổng tồn tại lâu; mất niềm tin người dùng.
\end{itemize}

Các tiêu chuẩn mở (AES, TLS) bền vững theo thời gian là vì hàng nghìn con mắt đã soi xét chúng, chứ không phải vì chúng bí mật, giấu kín.

\textbf{(c) (4 điểm) So sánh và đối chiếu giữa mật mã đối xứng và mật mã bất đối xứng:
    \begin{enumerate}
        \item Giải thích sự khác biệt cơ bản trong cách quản lý khóa của hai loại này.
        \item Đối với mỗi loại, hãy xác định các giả định toán học hoặc tính toán mà an ninh của chúng thường dựa vào.
        \item Mô tả một tình huống mà một loại sẽ rõ ràng là lựa chọn ưu việt hơn loại còn lại.
    \end{enumerate}
}

\begin{center}
\textbf{Trả lời:}
\end{center}

\begin{tabular}{|l|p{5cm}|p{5cm}|}
\hline
\textbf{Tiêu chí} & \textbf{Mật mã đối xứng} & \textbf{Mật mã bất đối xứng} \\
\hline
\textbf{Quản lý khóa} & 
1 khóa bí mật chung cho cả mã hóa và giải mã; cần phân phối khóa bí mật an toàn trước khi giao tiếp. & 
1 cặp khóa (công khai + riêng tư); khóa công khai chia sẻ rộng rãi, khóa riêng tư giữ bí mật; không cần kênh bí mật ban đầu. \\
\hline
\textbf{Cơ sở an ninh/giả định} & 
Dựa vào sức mạnh thuật toán và phân tích mật mã lâu dài (ví dụ: AES, HMAC-SHA-256). & 
Dựa vào các bài toán toán học khó (ví dụ: RSA dựa phân tích số nguyên lớn, DH/ECDH dựa logarit rời rạc). \\
\hline
\textbf{Tình huống ưu việt} & 
Mã hóa khối lượng dữ liệu lớn, yêu cầu hiệu năng cao (mã hóa đĩa, backup, streaming nội bộ sau khi đã có khóa phiên). & 
Trao đổi khóa ban đầu, xác thực/chữ ký số giữa các bên chưa có kênh bí mật (thiết lập TLS, ký cập nhật phần mềm). \\
\hline
\end{tabular}

\subsection{(10 điểm) Tính bí mật hoàn hảo (Perfect Secrecy)}

\textbf{(a) (3 điểm) Xét biến thể của one-time pad sử dụng phép AND bit ($\land$) thay vì XOR ($\oplus$):}

\begin{itemize}
    \item \textbf{Mã hóa:} $\text{Enc}(K, M) = K \land M$
    \item \textbf{Giải mã:} $\text{Dec}(K, C) = ?$
\end{itemize}

\begin{enumerate}
    \item \textbf{Scheme này có đúng không? Nếu đúng, hãy chỉ ra hàm giải mã. Nếu không, giải thích lý do.}
    
    \begin{center}
    \textbf{Trả lời:}
    \end{center}
    
    Scheme này \textbf{không đúng} (không đảm bảo tính đúng đắn).
    
    Phân tích chi tiết:
    \begin{itemize}
        \item Phép AND không khả nghịch vì có nhiều cặp $(K, M)$ khác nhau có thể tạo ra cùng ciphertext $C$.
        \item Ví dụ cụ thể: Nếu $C = 0100$, thì có thể có:
        \begin{itemize}
            \item $K = 0100, M = 0100 \rightarrow K \land M = 0100$
            \item $K = 0100, M = 0101 \rightarrow K \land M = 0100$
            \item $K = 0100, M = 0110 \rightarrow K \land M = 0100$
            \item $K = 0100, M = 0111 \rightarrow K \land M = 0100$
        \end{itemize}
        \item Khi giải mã với $K = 0100$ và $C = 0100$, không thể biết $M$ ban đầu là gì.
    \end{itemize}
    
    Tại sao AND không khả nghịch:
    \begin{itemize}
        \item Bit 0 trong $K$ ``che giấu'' thông tin về $M$ (luôn cho kết quả 0)
        \item Bit 1 trong $K$ ``tiết lộ'' thông tin về $M$ (kết quả = $M$)
        \item Thông tin bị mất không thể khôi phục được.
    \end{itemize}
    
    \item \textbf{Scheme này có đảm bảo bí mật hoàn hảo không? Giải thích.}
    
    \begin{center}
    \textbf{Trả lời:}
    \end{center}
    
    \textbf{Không đảm bảo bí mật hoàn hảo.}
    
    Phân tích chi tiết:
    \begin{itemize}
        \item Perfect Secrecy yêu cầu: $\Pr[\text{Enc}(K,M) = C] = 1/2^n$ với mọi $M$ và $C$
        \item Với phép AND, phân phối của ciphertext bị lệch nghiêm trọng
    \end{itemize}
    
    Ví dụ cụ thể với 2-bit:
    \begin{itemize}
        \item Nếu $M = 00$: chỉ có thể tạo ra $C = 00$ (dù $K$ là gì)
        \item Nếu $M = 01$: có thể tạo ra $C \in \{00, 01\}$ tùy theo $K$
        \item Nếu $M = 10$: có thể tạo ra $C \in \{00, 10\}$ tùy theo $K$
        \item Nếu $M = 11$: có thể tạo ra $C \in \{00, 01, 10, 11\}$ tùy theo $K$
    \end{itemize}
    
    Tấn công thông tin:
    \begin{itemize}
        \item Nếu attacker thấy $C = 11$, họ biết chắc $M = 11$ (vì chỉ có $M = 11$ mới tạo ra được $C = 11$)
        \item Nếu attacker thấy $C = 00$, họ biết $M$ có thể là bất kỳ giá trị nào (nhưng vẫn có thông tin)
        \item Điều này vi phạm định nghĩa Perfect Secrecy: attacker có thể phân biệt các $M$ khác nhau từ $C$.
    \end{itemize}
\end{enumerate}

\textbf{(b) (4 điểm) Xét biến thể one-time pad trên các chữ số thập phân (0-9):
    \begin{itemize}
        \item \textbf{Mã hóa:} $\text{Enc}(K, M) = (K + M) \bmod 10$
        \item \textbf{Giải mã:} $\text{Dec}(K, C) = (C - K) \bmod 10$
    \end{itemize}
}

Với $K, M, C \in \{0, 1, 2, \ldots, 9\}$

\begin{enumerate}
    \item \textbf{Chứng minh scheme này đúng}.
    
    \textbf{Chứng minh:}
    \begin{align}
        \text{Dec}(K, \text{Enc}(K, M)) &= ((K + M) \bmod 10 - K) \bmod 10 \\
        &= ((K + M - K) \bmod 10) \\
        &= (M \bmod 10) \\
        &= M \quad \text{(vì } M \in \{0,\ldots,9\}\text{)}
    \end{align}
    
    \textbf{Giải thích từng bước:}
    \begin{itemize}
        \item Bước 1: Thay thế $\text{Enc}(K,M) = (K + M) \bmod 10$
        \item Bước 2: Áp dụng tính chất phân phối của phép trừ trong modular arithmetic
        \item Bước 3: Đơn giản hóa $(K + M - K) = M$
        \item Bước 4: Vì $M \in \{0,1,\ldots,9\}$, nên $M \bmod 10 = M$
        \item Kết luận: $\text{Dec}(K, \text{Enc}(K,M)) = M$ với mọi $K$ và $M$
    \end{itemize}
    
    \item \textbf{Chứng minh scheme này đảm bảo bí mật hoàn hảo nếu $K$ chọn ngẫu nhiên đều.}
    
    \textbf{Chứng minh:}
    
    \textbf{Bước 1: Tính duy nhất của khóa}
    \begin{itemize}
        \item Với mỗi cặp $(M, C)$, tồn tại duy nhất một $K$ thỏa mãn: $C = (K + M) \bmod 10$
        \item Giải ra: $K = (C - M) \bmod 10$
    \end{itemize}
    
    \textbf{Bước 2: Phân phối đồng nhất của ciphertext}
    \begin{itemize}
        \item Nếu $K$ được chọn ngẫu nhiên đều từ $\{0,1,\ldots,9\}$, thì $\Pr[K = k] = 1/10$ với mọi $k$
        \item Với mỗi $M$ cố định, khi $K$ thay đổi từ 0 đến 9, ciphertext $C = (K + M) \bmod 10$ cũng chạy qua tất cả giá trị từ 0 đến 9
        \item Do đó: $\Pr[\text{Enc}(K,M) = C] = \Pr[K = (C-M) \bmod 10] = 1/10$ với mọi $M$ và $C$
    \end{itemize}
    
    \textbf{Bước 3: Kết luận Perfect Secrecy}
    \begin{itemize}
        \item Điều kiện Perfect Secrecy: $\Pr[\text{Enc}(K,M) = C] = 1/10$ với mọi $M$ và $C$
        \item Attacker không thể phân biệt các $M$ khác nhau từ $C$ vì phân phối của $C$ độc lập với $M$
    \end{itemize}
\end{enumerate}

Tổng kết: Cộng modulo là ``hoán vị dịch'' trên mỗi chữ số; nếu $K$ đồng nhất thì $K + M \bmod 10$ cũng đồng nhất, nên không lộ thông tin về $M$.

\textbf{(c) (3 điểm) Xét one-time pad với khóa ngắn bằng nửa độ dài thông điệp:}

\begin{itemize}
    \item $M = (M_1, M_2)$, với $|M_1| = |M_2| = |K|$
    \item $\text{Enc}(K, M) = (K \oplus M_1, K \oplus M_2)$
\end{itemize}

\textbf{Yêu cầu:} Đưa ra tấn công cụ thể phá vỡ tính bảo mật, chỉ rõ attacker có thể rút ra thông tin gì từ ciphertext.

\vspace{0.5em}
\textbf{Trả lời:}

\begin{enumerate}
    \item \textbf{Thu thập ciphertext:}
    \begin{itemize}
        \item Ciphertext thu được là $(C_1, C_2) = (K \oplus M_1, K \oplus M_2)$.
    \end{itemize}

    \item \textbf{Tấn công loại bỏ khóa:}
    \begin{itemize}
        \item Attacker tính $C_1 \oplus C_2$:
        \begin{align*}
            C_1 \oplus C_2 &= (K \oplus M_1) \oplus (K \oplus M_2) \\
            &= K \oplus M_1 \oplus K \oplus M_2 \\
            &= (K \oplus K) \oplus (M_1 \oplus M_2) \\
            &= 0 \oplus (M_1 \oplus M_2) \\
            &= M_1 \oplus M_2
        \end{align*}
        \item Như vậy, attacker biết được giá trị $M_1 \oplus M_2$.
    \end{itemize}

    \item \textbf{Phân tích thông tin bị lộ:}
    \begin{itemize}
        \item Giá trị $M_1 \oplus M_2$ tiết lộ mối quan hệ giữa hai nửa thông điệp.
        \item Ví dụ: Nếu $M_1 \oplus M_2 = 1010$, attacker có thể thử các giá trị $M_1$ để suy ra $M_2 = M_1 \oplus 1010$.
        \item Điều này giúp attacker loại trừ nhiều khả năng về thông điệp gốc, tức là ciphertext đã tiết lộ thông tin về plaintext.
    \end{itemize}

    \item \textbf{Giải thích vi phạm Perfect Secrecy:}
    \begin{itemize}
        \item Theo định nghĩa, Perfect Secrecy yêu cầu ciphertext không tiết lộ bất kỳ thông tin nào về plaintext.
        \item Tuy nhiên, ở đây attacker luôn biết $M_1 \oplus M_2$ từ ciphertext, nên scheme này không đạt Perfect Secrecy.
    \end{itemize}

    \item Việc sử dụng cùng một khóa $K$ cho hai phần thông điệp khiến khi XOR hai bản mã, khóa bị triệt tiêu, tương tự như bài toán two-time pad.

\end{enumerate}